\Needspace{15\baselineskip}
\section*{Des sections survivantes aux constantes fondamentales}
\addcontentsline{toc}{chapter}{Des sections survivantes aux constantes fondamentales}

Le continu généalogique déjà construit ne publie pas de simples développements décimaux : il publie des sections compatibles qui conservent orientation, frontière, retenue et mémoire. Une constante fondamentale est l'évaluation stable de l'une de ces sections sous un lecteur typé. La Partie III reconstruit donc l'objet, l'application et la preuve qui précèdent chaque valeur ; les comparaisons archimédiennes ou métrologiques n'apparaissent qu'après la sortie interne.

Au sens fixé par la
\hyperref[sec:tesis-inaugural-helicoidal-hmt-md]{tesis inaugural}, la lecture
nonadique n'agit pas à nouveau comme seuil et ne réinitialise pas le continu. La connexion nonadique conjointe et le prolongement
\[
w_6\longrightarrow w_{12}\longrightarrow w_{18}\longrightarrow w_{24}\longrightarrow w_{30}\longrightarrow R_{36}\longrightarrow G_9
\]
constituent, avec les changements et invariants fixés dans l'
\hyperref[sec:umbrales-prolongacion-tpk-inaugural]{architecture des seuils
TPK}, le premier tronçon certifié de la loi uniforme qui construit un unique
état coinductif enrichi. Parmi ses sorties corrélées, qui forment
le noyau minimal de son image nonadique typée, se trouvent
\(\pi_{\rm HMT},e_{\rm HMT},\varphi_{\rm HMT}\) et \(\alpha_{\rm HMT}\). L'unité d'origine n'aplanit pas leurs types : clôture angulaire, propagation, autoéchelle et couplage possèdent leurs propres lecteurs ; la coordonnée \(\alpha\) est publiée au moyen du registre signé, du sceau réversible et de la clôture dodécaphasique. Ce n'est qu'après la publication de ces sorties que s'ouvrent, avec des domaines séparés, la réalisation complexe, la mémoire modulaire, l'échelle d'action et la géométrie angulaire. L'état commun n'est pas épuisé par ces quatre coordonnées : ses lecteurs typés publient en outre des invariants, des agrégats, des relations spéculaires et des constantes mathématiques et physiques dérivées.
